\chapter{Assignment and Search}

\begin{proposition}
The dual assignment and the primal assignment problems yield the same solution.
\end{proposition}

\noindent \textbf{Proof}
The Primal problem is given by 
\begin{align*}
    \max_{\pi^{ih}} \quad & \sum_i \sum_h \pi^{ih} U^{ih} \\
    \text{s.t.} \quad & \sum_h \pi^{ih} = w^i, \quad \forall i \\
    & \sum_i \pi^{ih} = f^h, \quad \forall h 
\end{align*}

The Dual problem is given by 
\begin{align*}
    \min_{w^{i}, f^h} \quad & \sum_i  u^i w^i + \sum_h v^h f^h \\
    \text{s.t.} \quad & u^i + v^h \geq U^{ih}, \quad \forall i, h
\end{align*}

The Lagragean for the Primal problem is
\[
\mathcal{L}^P = \sum_i \sum_h  \pi^{ih} U^{ih} + \sum_h v^h \left( f^h - \sum_i \pi^{ih} \right) + \sum_i u^i \left(w^i - \sum_h \pi^{ih} \right)
\]

The Lagragean for the Dual problem is
\[
\mathcal{L}^D = \sum_i  u^i w^i + \sum_h v^h f^h + \sum_i \sum_h \pi^{ih} \left( U^{ih} -  u^i  - v^h \right)
\]

Obivously, when choosing proper Lagragean Multipliers,
namely, using $\{v^h\}, \{ u^i \}$ as the LM for the Primal problem and 
$\{\pi^{ih} \}$ for the Dual problem, 
these two Lagragean coincide with each other and thus yield the same solutions. \(\blacksquare\)


\begin{proposition}
    If $U^{ih}$ is supermodular, then the optimal matching is assortative.
\end{proposition}
\noindent \textbf{Proof}
Mathematically, it is a direct implication of the Hardy-Littlewood-Polya inequality, which says
\[
\langle x, y \rangle \leq \langle \hat{x}, \hat{y} \rangle,
\]
where $\hat{x}, \hat{y}$ are the decreasing rearrangements of $x, y$ respectively.

Now, we prove the proposition by contradiction with inductive steps.
Assume that the optimal mathching is not assortative;
concretely, assume that worker $i$ matches with firm $\sigma_*(i) \neq i$.
Then, using the supermodularity of $U$, we have 
\[
\sum_{i \neq 1 \neq {\sigma_*}^{-1}(1)} U^{i \sigma_*(i)} + U^{11} + U^{{\sigma_*}^{-1}(1) \sigma_*(1)} > \sum_i U^{i \sigma_*(i)}. 
\]
In the $k^{th}$ step, we have 
\[
\sum_{i \neq k \neq {\sigma_*}^{-1}(k)} U^{i \sigma_*(i)} + U^{kk} + U^{{\sigma_*}^{-1}(k) \sigma_*(k)} > \sum_i U^{i \sigma_*(i)}. 
\]
Iterate this procedure for all $i$, we reach a contradiction. \(\blacksquare\)


\section{Problems}
\subsection{Problem 3.1}
\paragraph{Problem.} Prove Proposition 3.2 (equivalence between stability of $(\pi, u, v)$ and the property that $\pi$ solves the primal assignment problem while $(u, v)$ solves its dual).

\paragraph{Solution.}
\textbf{``If'' side.} Suppose that $\pi$ is a solution to the primal problem and $(u, v)$ is a solution to the dual problem. By Lemma 2:
\[
U^{ih} = u^i + v^h.
\]
Since $\pi$ solves the primal problem:
\[
\sum_{i \in I} \sum_{h \in H} \pi^{ih} U^{ih}
= \sum_{i \in I} \sum_{h \in H} \pi^{ih} u^i + \sum_{h \in H} \sum_{i \in I} \pi^{ih} v^h
= \sum_{i \in I} w^i u^i + \sum_{h \in H} f^h v^h.
\]
Since $(u, v)$ solves the dual problem, for all $i \in I$ and $h \in H$ we have $u^i + v^h \geq U^{ih}$. Therefore, $(\pi, u, v)$ is a stable tuple.

\textbf{``Only if'' side.} Suppose that $(\pi_*, u_*, v_*)$ is a stable tuple. By absence of blocking pairs, for all $\pi \in \Pi$ we have:
\[
\sum_{i \in I} \sum_{h \in H} \pi^{ih} U^{ih} - \sum_{i \in I} \sum_{h \in H} \pi^{ih} u_*^i - \sum_{h \in H} \sum_{i \in I} \pi^{ih} v_*^h \leq 0.
\]
Therefore:
\[
\max_{\pi \in \Pi} \sum_{i \in I} \sum_{h \in H} \pi^{ih} U^{ih}
\leq \sum_{i \in I} w^i u_*^i + \sum_{h \in H} f^h v_*^h
\qquad \text{s.t.} \quad
\sum_{h \in H} \pi^{ih} = w^i, \quad \sum_{i \in I} \pi^{ih} = f^h.
\]
By feasibility:
\[
\sum_{i \in I} w^i u_*^i + \sum_{h \in H} f^h v_*^h
= \sum_{i \in I} \sum_{h \in H} \pi_*^{ih} U^{ih}
\leq \max_{\pi \in \Pi} \sum_{i \in I} \sum_{h \in H} \pi^{ih} U^{ih}.
\]
Therefore, $\pi_*$ solves the primal problem. Moreover,
\begin{align*}
\min_{(u,v) \in \mathbb{R}^I \times \mathbb{R}^H} \max_{\pi \in \Pi} S(\pi, u, v)
& \leq \max_{\pi \in \Pi} S(\pi, u_*, v_*) \\
& = S(0, u_*, v_*)
= \sum_{i \in I} w^i u_*^i + \sum_{h \in H} f^h v_*^h \\
& = \sum_{i \in I} \sum_{h \in H} \pi_*^{ih} U^{ih}
= \max_{\pi \in \Pi} \min_{(u,v) \in \mathbb{R}^I \times \mathbb{R}^H} S(\pi, u, v) \\
& = \min_{(u,v) \in \mathbb{R}^I \times \mathbb{R}^H} \max_{\pi \in \Pi} S(\pi, u, v).
\end{align*}
Therefore, $(u_*, v_*)$ solves the dual problem.


\subsection{Problem 3.2}
\paragraph{Problem.} Show that, if $I = H$ and $U$ is supermodular, the solution to the Monge problem is the identity permutation; that is, each agent/worker $i \in I$ is paired with firm $h = i$.

\paragraph{Solution.}
We prove by induction that an optimal solution $\sigma_*$ must be such that $\sigma_*(i) = i$ for all $i \in I$.
\begin{itemize}
\item \textit{Initial step.} Suppose that $\sigma_*$ is optimal but $\sigma_*(1) \neq 1$. Then:
\[
\sum_{1 \neq i \neq \sigma_*^{-1}(1)} U^{i \sigma_*(i)} + U^{11} + U^{\sigma_*^{-1}(1) \sigma_*(1)} > \sum_{i \in I} U^{i \sigma_*(i)}.
\]
Therefore, $\sigma_*$ is not an optimal solution.
\item \textit{Inductive hypothesis.} If $\sigma_*$ is an optimal solution, then $\sigma(i) = i$ for all $i < k$.
\item \textit{Inductive step.} Suppose that $\sigma_*$ is optimal but $\sigma_*(k) \neq k$. By inductive hypothesis, $\sigma_*(k) > k$ and $\sigma_*^{-1}(k) > k$. But then:
\[
\sum_{k \neq i \neq \sigma_*^{-1}(k)} U^{i \sigma_*(i)} + U^{kk} + U^{\sigma_*^{-1}(k) \sigma_*(k)} > \sum_{i \in I} U^{i \sigma_*(i)}.
\]
Therefore, $\sigma_*$ is not an optimal solution.
\end{itemize}